\vspace{0.5em}\noindent\textbf{Get the regional DNS name of S3 interface
endpoint}\par

\begin{enumerate}
\def\labelenumi{\arabic{enumi}.}
\item
  From the Amazon VPC menu, choose Endpoints.
\item
  Click the name of newly created endpoint: s3-interface-endpoint. Click
  details and save the regional DNS name of the endpoint (the first one)
  to your text-editor for later use.
\end{enumerate}

\begin{figure}
\centering
\pandocbounded{\includegraphics[keepaspectratio,alt={dns name}]{5-Workshop/5.4-S3-onprem/dns.png}}
\caption{dns name}
\end{figure}

\vspace{0.5em}\noindent\textbf{Connect to EC2 instance in ``VPC
On-prem''}\par

\begin{enumerate}
\def\labelenumi{\arabic{enumi}.}
\item
  Navigate to \textbf{Session manager} by typing ``session manager'' in
  the search box
\item
  Click \textbf{Start Session}, and select the EC2 instance named
  \textbf{Test-Interface-Endpoint}. This EC2 instance is running in
  ``VPC On-prem'' and will be used to test connectivty to Amazon S3
  through the Interface endpoint we just created. Session Manager will
  open a new browser tab with a shell prompt: \textbf{sh-4.2 \$}
\end{enumerate}

\begin{figure}
\centering
\pandocbounded{\includegraphics[keepaspectratio,alt={Start session}]{5-Workshop/5.4-S3-onprem/start-session.png}}
\caption{Start session}
\end{figure}

\begin{enumerate}
\def\labelenumi{\arabic{enumi}.}
\setcounter{enumi}{2}
\item
  Change to the ssm-user's home directory with command ``cd
  \textasciitilde{}''
\item
  Create a file named testfile2.xyz
\end{enumerate}

\begin{verbatim}
fallocate -l 1G testfile2.xyz
\end{verbatim}

\begin{figure}
\centering
\pandocbounded{\includegraphics[keepaspectratio,alt={user}]{5-Workshop/5.4-S3-onprem/cli1.png}}
\caption{user}
\end{figure}

\begin{enumerate}
\def\labelenumi{\arabic{enumi}.}
\setcounter{enumi}{4}
\tightlist
\item
  Copy file to the same S3 bucket we created in section 3.2
\end{enumerate}

\begin{verbatim}
aws s3 cp --endpoint-url https://bucket.<Regional-DNS-Name> testfile2.xyz s3://<your-bucket-name>
\end{verbatim}

\begin{itemize}
\tightlist
\item
  This command requires the --endpoint-url parameter, because you need
  to use the endpoint-specific DNS name to access S3 using an Interface
  endpoint.
\item
  Do not include the leading ' * ' when copying/pasting the regional DNS
  name.
\item
  Provide your S3 bucket name created earlier
\end{itemize}

\begin{figure}
\centering
\pandocbounded{\includegraphics[keepaspectratio,alt={copy file}]{5-Workshop/5.4-S3-onprem/cli2.png}}
\caption{copy file}
\end{figure}

Now the file has been added to your S3 bucket. Let check your S3 bucket
in the next step.

\vspace{0.5em}\noindent\textbf{Check Object in S3 bucket}\par

\begin{enumerate}
\def\labelenumi{\arabic{enumi}.}
\tightlist
\item
  Navigate to S3 console
\item
  Click Buckets
\item
  Click the name of your bucket and you will see testfile2.xyz has been
  added to your bucket
\end{enumerate}

\begin{figure}
\centering
\pandocbounded{\includegraphics[keepaspectratio,alt={check bucket}]{5-Workshop/5.4-S3-onprem/check-bucket.png}}
\caption{check bucket}
\end{figure}
